\documentclass[11pt,a4paper]{article}
\usepackage{fontspec,amsmath,amssymb,amsthm,microtype}
\setmainfont{lmroman10-regular.otf}[BoldFont=lmroman10-bold.otf,ItalicFont=lmroman10-italic.otf,BoldItalicFont=lmroman10-bolditalic.otf]
\usepackage{xcolor}
\usepackage[most]{tcolorbox}
\usepackage{mathrsfs}
\providecommand{\lt}{<}
\providecommand{\gt}{>}
\setmonofont{DejaVuSansMono.ttf}[BoldFont=DejaVuSansMono-Bold.ttf,ItalicFont=DejaVuSansMono-Oblique.ttf,BoldItalicFont=DejaVuSansMono-BoldOblique.ttf,Scale=0.86]
\newfontfamily\leanfallback{STIXTwoMath-Regular.otf}[Scale=0.86]
\definecolor{leanrule}{gray}{0.62}
\definecolor{leanback}{gray}{0.962}
\definecolor{leantitle}{gray}{0.22}
\definecolor{leanmuted}{gray}{0.38}
\newtcolorbox{leanbox}[1]{enhanced,breakable,colback=leanback,colframe=leanrule,
  boxrule=0.4pt,arc=0.6mm,left=2.4mm,right=2.4mm,top=1.8mm,bottom=1.8mm,
  before skip=0.8em,after skip=1em,fonttitle=\small\bfseries,coltitle=leantitle,
  colbacktitle=leanback,titlerule=0.3pt,title={#1}}
\newlength{\leancharwidth}
\newenvironment{leancode}{\par\smallskip\begingroup\ttfamily\footnotesize\raggedright
  \settowidth{\leancharwidth}{M}\setlength{\parskip}{0pt}\setlength{\parindent}{0pt}
  \frenchspacing}{\par\endgroup\smallskip}
\newcommand{\leanline}[2]{\hangindent=\dimexpr#1\leancharwidth+4\leancharwidth\relax
  \hangafter=1\noindent\hspace*{#1\leancharwidth}#2\par}
\newcommand{\leansource}[1]{\par{\footnotesize\color{leanmuted}#1}\par}
\newcommand{\leanlabel}[1]{\par\smallskip\noindent{\small\bfseries #1}\par}
\usepackage[margin=25mm]{geometry}
\usepackage[unicode,hidelinks]{hyperref}
\hypersetup{pdfauthor={LangevinLean Project}}
\setlength{\emergencystretch}{3em}
\setlength{\parskip}{0.4em}
\setlength{\parindent}{0pt}
\allowdisplaybreaks[2]
\raggedbottom
\title{\texorpdfstring{The exact anchored law of the actual random MALA path}{The exact anchored law of the actual random MALA path}}
\author{LangevinLean Project}
\date{}
\begin{document}\maketitle
Draft mathematical exposition; these notes may contain mistakes.


The finite-time pair of the actual initial state and sampled MALA endpoint has law \(\pi\otimes K^n\), with the explicit fresh-noise kernel. This identifies actual random endpoint pairs with the proved reversible kernel powers.


\section{\texorpdfstring{Induction on the actual anchored law}{Induction on the actual anchored law}}


\textbf{Claim.}

\[\operatorname{Law}(X_0,X_n^{\mathrm{MALA}})=\pi\otimes K^n,\qquad \operatorname{Law}(X_0,X_n^{\mathrm{MALA}})=\operatorname{Law}(X_n^{\mathrm{MALA}},X_0).\]

\textbf{Proof.} At time zero the anchor/\allowbreak{}endpoint law is the diagonal Gibbs law, exactly \(\pi\otimes K^0\). Both the sampled anchored recursion and the explicit kernel-power pair advance by the same parallel-kernel update. Induction gives equality of the full pair laws. The actual product initial/\allowbreak{}innovation law transfers this equality to the canonical sampled path. Reversal follows from the proved explicit-power detailed balance.

\begin{leanbox}{Lean: \texttt{anchored\_\allowbreak{}chainLaw}, \texttt{recursion\_\allowbreak{}law}, \texttt{anchored\_\allowbreak{}state\_\allowbreak{}law} and 1 more}
\leansource{\texttt{anchored\_\allowbreak{}chainLaw} \quad WassersteinMixing/\allowbreak{}SemigroupAnchoredMALALaw.lean, lines 23–39}
\begin{leancode}
\leanline{0}{theorem\ anchored\_chainLaw\ (V\ :\ E\ →\ ℝ)\ (h\ :\ ℝ)\ (L\ :\ ℝ≥0)}
\leanline{4}{(hdiff\ :\ Differentiable\ ℝ\ V)\ (hLip\ :\ LipschitzWith\ L\ (gradient\ V))}
\leanline{4}{(μ\ :\ Measure\ E)\ [IsProbabilityMeasure\ μ]\ (n\ :\ ℕ)\ :}
\leanline{4}{ULA.chainLaw\ (anchoredUpdate\ V\ h)\ (μ.map\ (fun\ x\ ↦\ (x,\ x)))}
\leanline{6}{(Measure.infinitePi\ (fun\ \_\ :\ ℕ\ ↦\ innovationLaw))\ n\ =}
\leanline{8}{μ\ ⊗ₘ\ StationaryKernelPower.kernelPower\ (MALAStationaryPair.explicitKernel\ V\ h)\ n}
\end{leancode}
\leansource{\texttt{recursion\_\allowbreak{}law} \quad WassersteinMixing/\allowbreak{}SemigroupAnchoredMALALaw.lean, lines 48–70}
\begin{leancode}
\leanline{0}{theorem\ recursion\_law\ \{H\ :\ Type*\}\ [MeasurableSpace\ H]}
\leanline{4}{(F\ :\ H\ →\ ((E\ P.dimension)\ ×\ I)\ →\ H)\ (hF\ :\ Measurable\ (Function.uncurry\ F))}
\leanline{4}{(j\ :\ E\ P.dimension\ →\ H)\ (hj\ :\ Measurable\ j)}
\leanline{4}{(r\ :\ ℝ)\ (hr\ :\ 0\ <\ r)\ (n\ :\ ℕ)\ :}
\leanline{4}{(probability\ P\ V).map\ (fun\ ω\ ↦\ ULA.noiseDrivenChain\ F}
\leanline{6}{(j\ (input\ P\ V\ hV\ r\ ω).1)\ (input\ P\ V\ hV\ r\ ω).2\ n)\ =}
\leanline{6}{ULA.chainLaw\ F\ ((ULA.gibbsMeasure\ V).map\ j)}
\leanline{8}{(Measure.infinitePi\ (fun\ \_\ :\ ℕ\ ↦\ innovationLaw))\ n}
\end{leancode}
\leansource{\texttt{anchored\_\allowbreak{}state\_\allowbreak{}law} \quad WassersteinMixing/\allowbreak{}SemigroupAnchoredMALALaw.lean, lines 73–86}
\begin{leancode}
\leanline{0}{theorem\ anchored\_state\_law\ (r\ :\ ℝ)\ (hr\ :\ 0\ <\ r)\ (n\ :\ ℕ)\ :}
\leanline{4}{(probability\ P\ V).map}
\leanline{6}{(fun\ ω\ ↦\ (actualInitial\ P\ V\ hV\ ω.1,\ (state\ P\ V\ hV\ r\ n\ ω).2))\ =}
\leanline{8}{MALAStationaryPair.explicitPowerJointLaw\ V\ r\ n}
\end{leancode}
\leansource{\texttt{anchored\_\allowbreak{}state\_\allowbreak{}swap} \quad WassersteinMixing/\allowbreak{}SemigroupAnchoredMALALaw.lean, lines 89–97}
\begin{leancode}
\leanline{0}{theorem\ anchored\_state\_swap\ (r\ :\ ℝ)\ (hr\ :\ 0\ <\ r)\ (n\ :\ ℕ)\ :}
\leanline{4}{((probability\ P\ V).map}
\leanline{6}{(fun\ ω\ ↦\ (actualInitial\ P\ V\ hV\ ω.1,\ (state\ P\ V\ hV\ r\ n\ ω).2))).map\ Prod.swap\ =}
\leanline{6}{(probability\ P\ V).map}
\leanline{8}{(fun\ ω\ ↦\ (actualInitial\ P\ V\ hV\ ω.1,\ (state\ P\ V\ hV\ r\ n\ ω).2))}
\end{leancode}
\leanlabel{How the Lean proof follows the argument}
\texttt{anchored\_\allowbreak{}chain\allowbreak{}Law} compares the zero-step and successor laws. \texttt{recursion\_\allowbreak{}law} applies the actual sampled-input law. \texttt{anchored\_\allowbreak{}state\_\allowbreak{}law} identifies the true canonical pair and \texttt{anchored\_\allowbreak{}state\_\allowbreak{}swap} substitutes the proved explicit-power swap identity.
\end{leanbox}


Earlier results used here: \cite[Section~3.20.1]{LangevinLean2026ULASC}; \cite[Section~3.20.4]{LangevinLean2026ULASC}; \cite[Section~3.20.5]{LangevinLean2026ULASC}; \cite[Section~3.26.3]{LangevinLean2026ULASC}; \cite[Section~3.33.6]{LangevinLean2026ULASC}; \cite[Section~3.33.9]{LangevinLean2026ULASC}.

\begin{thebibliography}{ULA-LSI}
\bibitem[ULA-SC]{LangevinLean2026ULASC}
LangevinLean Project. \emph{Unadjusted Langevin under Strong Convexity}. LL-01, 2026. Draft report, content revision d34ead06987c. \url{https://odakyildiz.github.io/LangevinLean/docs/pdf/UnadjustedLangevin/FullProof.pdf}.
\end{thebibliography}

\end{document}
